\scanpage{214}
\mspage{69}

5. \uline{Dualité formelle}

Soient $P,Q$ deux ensembles ordonnés. \deleted{\unsure{Une|Nous}} J'\unsure{essaye|envisage|essaie} de dégager \deleted{la} \inserted{une} notion de ``dualité'' (formelle) \deleted{entre} \inserted{reliant} $P$ et $Q$. Ce sera une relation entre $P$ et $Q$

\[
R\subset P\times Q
\tag*{(5.1.)}
\]

ayant les propriétés suivantes

(D) Pour toute partie $A$ de $P$, l'ensemble

\[
\begin{aligned}
 A^{\unsure{\cup|\smile}}\overset{\mathrm{def}}{=}
 \big\{q\in Q\mid{}&\deleted{\unsure{R}}(p,q)\in R\ \deleted{\unsure{pour}\ \unsure{A}}\ \text{pour tout }p\in A\\
 &\text{i.e.}\quad A\times\{q\}\subset R\qquad\big\}
\end{aligned}
\]

est une partie ouverte de $Q$ ayant \inserted{un} plus grand élément $\beta(A)$\unsure{:|.} \inserted{$A^\cup=Q_{/\beta(A)}$.} Dualement, pour \deleted{toute} partie $B$ de $Q$, l'ensemble

\[
\begin{aligned}
 B^{\unsure{\cap|\frown}}\overset{\mathrm{def}}{=}
 \big\{p\in P\mid{}&(p,q)\in R\ \text{pour }\unsure{\text{tout}}\ q\in B,\\
 &\text{i.e.}\qquad \{p\}\times B\subset R\qquad\big\}
\end{aligned}
\]

est une partie ouverte de $P$ ayant un plus grand élément $\alpha(B)$, i.e. $B^\cap=P_{/\alpha(B)}$.

Le plus souvent, $P$ et $Q$ seront stables par sup quelconques (donc aussi par inf quelconques). Dans ce cas, (D) équivaut, du côté des $A^\cup$, (pour $A\in\unsure{\mathfrak P|\mathcal P}(P)$) à ceci

$(D_{0\unsure{a|\bullet|*}})$ Pour tout $p\in P$, $\deleted{\unsure{R(p)=\{q\in}}$ $R(p)=\{q\in Q\mid(p,q)\in R\}$ est une partie ouverte \uline{filtrante} de $\deleted{\unsure{P|R}}Q$, ou ce qui revient au $=$, une partie non vide ($\Longleftrightarrow$ contenant le plus petit élément $\varnothing_Q$ de $Q$) stable par sup de deux éléments. De plus $R(p)$ est stable par \deleted{\unsure{lim|les}} Sup filtrants, ou ce qui revient au $=$ (étant lui même filtrant), il contient un plus grand élément. Il revient
% archive sheet number: 214
